% Generated draft: requires mathematical and source review.
\documentclass[11pt]{article}
\usepackage[margin=25mm]{geometry}
\usepackage{fontspec}
\setmainfont{Latin Modern Roman}
\usepackage{amsmath,amssymb,mathtools,unicode-math}
\setmathfont{Latin Modern Math}
\usepackage{enumitem,needspace,xurl,hyperref}
\setlength{\parindent}{0pt}
\setlength{\parskip}{5pt}
\setlength{\emergencystretch}{3em}
\newcommand{\N}{\mathbb{N}}
\newcommand{\Z}{\mathbb{Z}}
\newcommand{\Q}{\mathbb{Q}}
\newcommand{\R}{\mathbb{R}}
\newcommand{\C}{\mathbb{C}}
\newcommand{\F}{\mathbb{F}}
\newcommand{\A}{\mathbb{A}}
\newcommand{\PP}{\mathbb{P}}
\newcommand{\E}{\mathbb{E}}
\newcommand{\eps}{\varepsilon}
\newcommand{\dd}{\,\mathrm{d}}
\begin{document}
\section*{UnsolvedMath ID: 30004795}
\textbf{Source problem code:} catalogue entry
\subsection{Definitions and mathematical statement}
Let \(k\) be a non-Archimedean local field and let \(G\) be a reductive algebraic group over \(k\); write \(G(k)\) for its group of \(k\)-rational points. Let \(\pi\) be an essentially discrete-series representation of \(G(k)\), and let \(\phi\) be its Langlands parameter, in the formal-degree setting intended by the source report.

The supplied record uses \(\operatorname{Ad}_G\) for the adjoint representation occurring in that setting, \(\omega\) for the orthogonal root-number invariant, \(\phi_{\operatorname{prin}}\) for the principal parameter, \(z_{\operatorname{Ad}_G}\) for the associated central involution, and \(\chi_\phi\) for the central character attached to \(\phi\). These descriptions retain the source notation; their precise constructions and conventions are not supplied.

\textbf{Conjectural identity (reported as Conjecture 1.2).} For every such \(k\), \(G\), and \(\phi\), subject to the hypotheses and conventions of the report's preceding formal-degree setting, prove that
\[
\frac{\omega\bigl(\operatorname{Ad}_G\circ\phi\bigr)}
     {\omega\bigl(\operatorname{Ad}_G\circ\phi_{\operatorname{prin}}\bigr)}
=\chi_\phi\bigl(z_{\operatorname{Ad}_G}\bigr).
\]

\textbf{Unresolved qualifications.} The original statement explicitly refers to a preceding formal-degree setting that is absent from the supplied record. In particular, the record does not specify the precise parameter category, the adjoint representation and its underlying space, the definition and normalization of \(\omega\), the construction of the principal parameter, or the group containing \(z_{\operatorname{Ad}_G}\) on which \(\chi_\phi\) is evaluated. It also does not state whether additional restrictions on \(k\), \(G\), or \(\pi\) are imposed there. Consequently, the display preserves the supplied conjecture but does not supply a complete standalone formulation without recovering those conventions and hypotheses.
\subsection{Short English statement}
The conjecture equates a ratio of orthogonal root numbers for a discrete-series Langlands parameter and the principal parameter with a central-character value at an associated central involution. The source's preceding formal-degree setting is needed to define these quantities precisely. The supplied partially solved label and literature assessment are dated curation, not independently verified conclusions about this identity.
\subsection{Sources}
\begin{itemize}
\item Supplied local registry record: permanent UnsolvedMath ID \(30004795\), title \emph{Root Numbers and Central Characters of Langlands Parameters}, exact source problem code \texttt{catalogue entry}, in the Oberwolfach Reports Open Problems collection. The proposed year is recorded as 2021; no proposer is named.
\item Supplied report citation: \emph{Automorphic Forms, Geometry and Arithmetic (2022)}, Oberwolfach Reports, pp.\ 2089--2156, DOI \texttt{10.4171/owr/2021/39}. \url{https://doi.org/10.4171/owr/2021/39}. The registry attributes the identity to Conjecture 1.2 and records that its curator inspected the report and restored the surrounding local-field, reductive-group, representation, and principal-parameter context. This verification claim is supplied provenance, not an inspection performed for this draft.
\item Supplied literature source, labelled ``original Oberwolfach report'': \url{https://doi.org/10.4171/owr/2021/39}.
\item Supplied literature source, labelled ``Beuzart-Plessis, classical groups'': \url{https://arxiv.org/abs/2508.08470}.
\item Supplied literature source, labelled ``formal degree of regular supercuspidals'': \url{https://doi.org/10.4171/JEMS/1412}.
\item Literature assessment date: \texttt{2026-08-22}, copied exactly from the supplied record. The supplied status is \texttt{partially\_solved}. Its assessment states that the formal-degree formula is proved for broad families, including classical groups, regular Yu supercuspidals, and principal-series discrete series, but is not established for all reductive \(p\)-adic groups and discrete-series parameters. This is supplied dated curation derived from the OWR material, not independently verified literature assessment; its time-sensitive status and cited claims require independent checking. The supplied assessment does not explain how these formal-degree results apply to the displayed root-number identity.
\end{itemize}
\end{document}
